
\section{Introdução}


\noindent \begin{minipage}[c]{0.6\textwidth}
\par O uso de banco e dados teve seu início com o advento da computação, para o armazenar e gerenciar dados, com as necessidades aumentando diversos pesquisadores foram desenvolvendo e aprimorando uma marco importante segundo \cite{histdb:2024}, em meados de 1980 a IBM desenvolve a \textit{Structured Query Language} (SQL), com o passar dos tempo o ícone que represenda um banco de dados é o da figura \ref{fig:icodb}. Seguindo os avanços e nesseciades tecnilogicas temos dois grandes tipos de banco de dados: os \textbf{Relacionais} e os \textbf{Não relacionais}, dos quais o objeto tema desta aula prática seráo \textbf{Não Relacionais}.
\par Esta aula pratica tem por objetivo implementação de um banco de dados \textbf{Não relacional} em um contexto de um loja, a referida atividade será desenvolvida em duas etapas: Primeira atividade \ref{1atividade} e Segunda atividade \ref{2atividae}


\end{minipage}
\begin{minipage}[c]{0.4\textwidth}
  \captionof{figure}{Ícone de banco de dados}
  \includegraphics[width=\textwidth]{figure/pngegg.png}
  \label{fig:icodb}

  {\fontsize{10pt}{\baselineskip}\selectfont
    Fonte: \citeonline{DB:2024}}
\end{minipage}


\section{Métodos}

\par É porposto a elaboração de uma aula prática, tendo como fim a aplicação dos conhecimentos adqueridos em sala de aula virtual, e para tal é disponibilizado um \href{https://github.com/ENGENHARIA-DE-SOFTWARE-UNOPAR/NoSQL/blob/main/Roteiro-aula%20pratica.pdf}{roteiro aula prática}. Utiliza-se o \textit{software} \textbf{MongoDB} \ref{fig:logodb}, o referido \textit{software} é de código aberto e seu código pode ser consultado no \href{https://github.com/mongodb}{repositório do GitHub}.

\begin{figure}[H]
  \caption{Logo MongoDB}
  \includegraphics[scale=0.8]{figure/mongologo.png}
  \label{fig:logodb}

  {\fontsize{10pt}{\baselineskip}\selectfont
    Fonte: \citeonline{mongo:2024}}
\end{figure}




\section{Resultados}
\subsection{Primeira atividade}\label{1atividade}
\par Para a primeira atividade é proposto a criação de um banco de dados no \textit{softwarwe} MongoDB Compass entitulado de \textbf{lojadb}, esta banco de dados será utilizado em um contexto de uma loja o qual deverá cadastrar documentos referentes a vendas efetuadas pela referida loja, e subdivide-se em duas etapas, a criação do banco de dados conforme tabela \ref{tab:banco}, e inserção de dados, no caso, endereço dos clientes conforme tabela \ref{tab:end} e produtos adiqiridos pelos cliente conforme listagem \ref{tab:com}.

\begin{table}[H]
\center
  \caption{Criação do banco de dados.}
  \begin{tabular}{ | c | c  | c | c |}
      \hline
      \textbf{Nome} & \textbf{Cliente VIP} & EMAIL & TELEFONE \\
          & (1-SIM / 0-Não)  &  &    \\
      \hline
          João & 1 & joao@email.com & 9999-1111, 8888-1111 \\
      \hline
          Marcos & 0 &  & 9999-2222 \\
      \hline
          Maria & 1 &maria@email.com & 9999-3333, 8888-3333, 9988-3000 \\
      \hline
  \end{tabular}
  \label{tab:banco}
  \flushleft % Alinha a esquerda a referencia
  %\raggedleft

  {\fontsize{10pt}{\baselineskip}\selectfont
    Fonte: \citeonline{aula:2024}}
\end{table}

\par Para a implementação da tabela \ref{tab:banco}, foi confecionado o código \ref{cod:banco}.


\begin{table}[H]
  \caption{Endereços a serem adicionados para os clientes.}
  \center
  \begin{tabular}{ | c | c |}
      \hline
      \textbf{Nome} & \textbf{Endereço} \\
      \hline
          João & Rua Um, 1000, Apto 1 Bloco 1. São Paulo/SP \\
      \hline
          Marcos & Rua Dois, 4000. Campinas/SP \\
      \hline
          Maria & Rua Três, 3000. Londrina/PR \\
      \hline
  \end{tabular}
  \label{tab:end}
  \flushleft % Alinha a esquerda a referencia
  %\raggedleft

  {\fontsize{10pt}{\baselineskip}\selectfont
    Fonte: \citeonline{aula:2024}}
\end{table}

\par Para a inserçao dos dados da tebela \ref{tab:end}, foi confecionado a listagem \ref{cod:end}.

\par Para concluir esta etapa e utilizar varias formas do MongoDB, sugere a inserção de mais dados como demonstrado na tabela \ref{tab:com}, inserindo os produtos adqueridos pelos clientes, como: Nome do equipamento, preço e quantidade.

\begin{table}[H]
  \caption{Produtos adquiridos pelos clientes}
  \center
  \begin{tabular}{ | c | c |}
      \hline
      \textbf{Nome} & \textbf{Compras} \\
      \hline
           &  \multicolumn{1}{l|}{Nome do Produto: caderno} \\
        João  & \multicolumn{1}{l|}{Preço: R\$: 5.000,00}  \\
          & \multicolumn{1}{l|}{Quantidade: 1 } \\
      \hline
           &  \multicolumn{1}{l|}{Nome do Produto: caderno} \\
          & \multicolumn{1}{l|}{Preço: R\$: 20,00}  \\
          & \multicolumn{1}{l|}{Quantidade: 1 } \\
        Marcos  & \\
          & \multicolumn{1}{l|}{Nome do Produto: Borracha} \\
          & \multicolumn{1}{l|}{Preço: R\$: 2,00} \\
          & \multicolumn{1}{l|}{Quantidade: 2} \\
      \hline
           &  \multicolumn{1}{l|}{Nome do Produto: Tablet} \\
          & \multicolumn{1}{l|}{Preço: R\$: 2.500,00}  \\
          & \multicolumn{1}{l|}{Quantidade: 1 } \\
        Maria  & \\
          & \multicolumn{1}{l|}{Nome do Produto: Capa para tablet} \\
          & \multicolumn{1}{l|}{Preço: R\$: 50,00} \\
          & \multicolumn{1}{l|}{Quantidade: 1} \\
      \hline

  \end{tabular}
  \label{tab:com}
  \flushleft % Alinha a esquerda a referencia
  %\raggedleft

  {\fontsize{10pt}{\baselineskip}\selectfont
    Fonte: \citeonline{aula:2024}}
\end{table}

\par Para a realização desta atividades disponibilizo o \textit{scripty} \ref{cod:com} da tabela \ref{tab:com}.




\subsection{Segunda Atividade}\label{2atividae}

%\lstinputlisting[caption={consulta.sql}, label={cod:com}]{consulta.sql}
%\lstinputlisting[language=c, caption={código externo}, label={cod:externo}, captionpos=t]{main.c} %Busca os codigos na pasta /cod

%%%%%%%%%%%%%%%%%%%%%%%%%%%%%%%%%%%%%%%
%\begin{figure}[H] %Figuras da aula pratica 1.1
%  \center
%  \caption{Resultado da atividade prática 1.2,
%  \subfigure[ Algoritmo.\label{fig:pri2}]{\includegraphics[scale=0.4]{figure/placeholder.jpg}}
%  \subfigure[Comportamento.\label{fig:seg2}]{\includegraphics[scale=.4]{figure/placeholder.jpg}}

%{\fontsize{10pt}{\baselineskip}\selectfont  Fonte: \citeonline{aula:2024}}
%\end{figure}

%%%%%%%%%%%%%%%%%%%%%%%%%%%%%%%%%%%%%%%%%%%%%%%%%%%%%%%%%%


Para referenciar utilize \cite{ninguem2022curioso}. Também pode ser citado integrada ao texto, de acordo com \citeonline{alguem2022nada}.

Para inserir imagens adicione a figura no diretório \textit{/figure}






\begin{enumerate}[label=\Roman{*}, ref=(\roman{*})]
  \item fsfsdf
  \item kugfhiuh
\end{enumerate}

\begin{asparaenum}
\item Anterior ... \cite{ninguem2022curioso}
\item Próximo ... \label{pl1}
\end{asparaenum}





\section{Conclusões}



  %$X \xLongleftarrow[\text{NATAN}]{\text{OGLIARI}} Y $ %COM TEXTO
	% $\uparrow$ %Seta para Cima
	%$\overleftarrow{NATAN}$
